\section{Mixed antiderivatives}
\label{sec:antiderivatives}

The antiderivatives below are functions of two lengths, which are written as $a$ and $b$ like in the thesis appendix.
This does not clash with the cell sizes: the touching moments evaluate them only at the points $(a,b)$, $(2a,b)$,
$(a,2b)$ and $(2a,2b)$ (Section~\ref{sec:touching}).

\defTitle{Mixed antiderivatives}{antiderivative} Let $Q=[0,\infty)^2$ be the closed first
quadrant and $Q^\circ$ its interior. For $p>-2$, $m,n\in\N_0$ and $(a,b)\in Q$ put
\begin{equation}
    \label{eq:G}
    \Gf{p}{m}{n}(a,b) = \int_0^a\!\!\int_0^b u^m v^n\bigl(u^2+v^2\bigr)^{p/2}\D v\D u .
\end{equation}
For a function $F$ on $Q$ and a rectangle $A\times B\subset Q$ with $A=[a_0,a_1]$ and $B=[b_0,b_1]$, the corner
difference of $F$ over $A\times B$ is
\begin{equation}
    \cornerdiff_{A\times B}F = F(a_1,b_1)-F(a_0,b_1)-F(a_1,b_0)+F(a_0,b_0).
\end{equation}

\begin{figure}[ht]
    \centering
    % The corner difference (Proposition gproperties): a rectangle of Q as an alternating sum of the rectangles
% [0,U] x [0,V] at its four corners; on a piece that touches the singular point three corners lie on the axes.
\begin{tikzpicture}[line cap=round, line join=round, scale=1.15]
    \foreach \sh/\uo/\ui/\vo/\vi in {0/1/2.2/0.7/1.6, 4.2/0/1.6/0/1.2} {
        \begin{scope}[shift={(\sh,0)}]
            \begin{scope}
                \clip (0,0) rectangle (2.6,2.1);
                \foreach \r in {0.4,0.8,...,3.3} { \draw[PairedE, thin] (0,0) circle (\r); }
            \end{scope}
            \fill[PairedA] (\uo,\vo) rectangle (\ui,\vi);
            \draw[dashed] (0,\vi) -- (\ui,\vi) -- (\ui,0);
            \draw[dashed] (0,\vo) -- (\ui,\vo);
            \draw[dashed] (\uo,0) -- (\uo,\vi);
            \draw[thick] (\uo,\vo) rectangle (\ui,\vi);
            \draw[thick,->] (0,0) -- (2.8,0) node[anchor=west]{$u$};
            \draw[thick,->] (0,0) -- (0,2.3) node[anchor=south]{$v$};
            \fill[PairedF] (0,0) circle (2pt);
        \end{scope}
    }
    \node[anchor=north] at (1,0) {$a_0$};
    \node[anchor=north] at (2.2,0) {$a_1$};
    \node[anchor=east] at (0,0.7) {$b_0$};
    \node[anchor=east] at (0,1.6) {$b_1$};
    \foreach \u/\v/\s in {2.2/1.6/+, 1/0.7/+, 1/1.6/-, 2.2/0.7/-} {
        \fill (\u,\v) circle (1.6pt);
        \node[anchor=south west, inner sep=1.5pt, font=\small] at (\u,\v) {$\s$};
    }
    \begin{scope}[shift={(4.2,0)}]
        \node[anchor=north] at (1.6,0) {$a$};
        \node[anchor=east] at (0,1.2) {$b$};
        \fill (1.6,1.2) circle (1.6pt);
        \node[anchor=south west, inner sep=1.5pt, font=\small] at (1.6,1.2) {$+$};
        \foreach \u/\v in {0/1.2, 1.6/0} { \draw[fill=white] (\u,\v) circle (1.6pt); }
            \end{scope}
\end{tikzpicture}

    \caption{Left: the corner difference~\eqref{eq:corner}. The integral over the shaded rectangle is the alternating
    sum of $\Gf{p}{m}{n}$ at its corners, each value being the integral over $[0,a]\times[0,b]$ (dashed). The circles
    are level lines of $\sqrt{u^2+v^2}$ around the singular point (red). Right: a piece that touches the singular
    point. Three of its corners lie on the axes, where $\Gf{p}{m}{n}$ vanishes, and only $\Gf{p}{m}{n}(a,b)$ remains.}
    \label{fig:corner}
\end{figure}

\propTitle{Properties of $G$}{gproperties} Let $p>-2$ and $m,n\in\N_0$.
\begin{enumerate}[(i)]
    \item $\Gf{p}{m}{n}$ is continuous on $Q$ and vanishes on both axes: $\Gf{p}{m}{n}(a,0)=\Gf{p}{m}{n}(0,b)=0$.
    \item For every rectangle $[a_0,a_1]\times[b_0,b_1]\subset Q$,
    \begin{equation}
        \label{eq:corner}
        \int_{a_0}^{a_1}\!\!\int_{b_0}^{b_1} u^mv^n\bigl(u^2+v^2\bigr)^{p/2}\D v\D u
        = \cornerdiff_{[a_0,a_1]\times[b_0,b_1]}\Gf{p}{m}{n} .
    \end{equation}
    \item $\partial_a\partial_b\Gf{p}{m}{n}=a^mb^n\bigl(a^2+b^2\bigr)^{p/2}$ on $Q^\circ$, and $\Gf{p}{m}{n}$ is the only
    function in $C(Q)\cap C^2(Q^\circ)$ with this property that vanishes on both axes.
    \item $\Gf{p}{m}{n}(\lambda a,\lambda b)=\lambda^{m+n+p+2}\,\Gf{p}{m}{n}(a,b)$ for $\lambda>0$.
\end{enumerate}

\vspace{\baselineskip}
\proof
(i) The integrand is locally integrable on $Q$ because $p>-2$; continuity follows by dominated convergence, and the
integral over a rectangle with a side of length zero is zero.

(ii) This is inclusion--exclusion (Figure~\ref{fig:corner}): the rectangle $[a_0,a_1]\times[b_0,b_1]$ is
$[0,a_1]\times[0,b_1]$ minus $[0,a_0]\times[0,b_1]$ and $[0,a_1]\times[0,b_0]$, plus their overlap
$[0,a_0]\times[0,b_0]$.

(iii) The derivative follows from the fundamental theorem of calculus in each variable. If $D$ is the difference of two
such functions, $\partial_a\partial_bD=0$ gives $D(a,b)=e(a)+c(b)$ on $Q^\circ$; letting $b\to0^+$ and $a\to0^+$ with
$D=0$ on the axes shows that $e$ and $c$ are constants with $e+c=D(0,0)=0$.

(iv) Substitute $u=\lambda u'$ and $v=\lambda v'$ in~\eqref{eq:G}. \qed

\lemmaTitle{One-dimensional integrals}{oned} Let $a\ge0$, $b>0$, $m\in\N_0$ and $q\in\{-1,1,3,5\}$, and let $R=\sqrt{a^2+b^2}$. The integrals
\begin{equation}
    K^{(q)}_m(a)=\int_0^a u^m\bigl(u^2+b^2\bigr)^{q/2}\D u
\end{equation}
are
\begin{align}
    K^{(q)}_0 &= \begin{dcases}
        \arsinh\frac{a}{b}, & q=-1,\\[2pt]
        \frac12\Bigl(aR+b^2\arsinh\frac{a}{b}\Bigr), & q=1,\\[2pt]
        \frac18\Bigl(2aR^3+3b^2aR+3b^4\arsinh\frac{a}{b}\Bigr), & q=3,\\[2pt]
        \frac{1}{48}\Bigl(8aR^5+10b^2aR^3+15b^4aR+15b^6\arsinh\frac{a}{b}\Bigr), & q=5,
    \end{dcases}\label{eq:K0}\\
    K^{(q)}_1 &= \frac{R^{q+2}-b^{q+2}}{q+2}, \label{eq:K1}\\
    K^{(q)}_m &= \frac{a^{m-1}R^{q+2}-(m-1)b^2K^{(q)}_{m-2}}{m+q+1}, \qquad m\ge2, \label{eq:Kk}
\end{align}
and, for the logarithmic terms produced by the inner integrations,
\begin{equation}
    \label{eq:ibp}
    \int_0^a u^m\arsinh\frac{b}{u}\D u = \frac{1}{m+1}\Bigl(a^{m+1}\arsinh\frac{b}{a} + b\,K^{(-1)}_m(a)\Bigr).
\end{equation}
The same formulas with $a$ and $b$ exchanged give the inner integrals $\int_0^b v^n(u^2+v^2)^{p/2}\D v$.

\vspace{\baselineskip}
\proof
The base cases follow from
\begin{equation}
    \begin{aligned}
    \frac{\D}{\D u}\arsinh\frac{u}{b} &= \bigl(u^2+b^2\bigr)^{-1/2},\qquad
    \frac{\D}{\D u}\bigl(u^2+b^2\bigr)^{(q+2)/2} = (q+2)\,u\,\bigl(u^2+b^2\bigr)^{q/2},\\
    \frac{\D}{\D u}\Bigl(u\bigl(u^2+b^2\bigr)^{q/2}\Bigr) &= (q+1)\bigl(u^2+b^2\bigr)^{q/2}-qb^2\bigl(u^2+b^2\bigr)^{(q-2)/2};
    \end{aligned}
\end{equation}
the first one gives $K^{(-1)}_0$ and the second one~\eqref{eq:K1}. The last one gives the recursion
\begin{equation}
    K^{(q)}_0=\frac{aR^q+qb^2K^{(q-2)}_0}{q+1},
\end{equation}
which produces~\eqref{eq:K0} from $K^{(-1)}_0$.

For the recursion with $m\ge2$, integrate from $0$ to $a$
\begin{equation}
    \frac{\D}{\D u}\bigl(u^{m-1}(u^2+b^2)^{(q+2)/2}\bigr)=(m-1)b^2u^{m-2}(u^2+b^2)^{q/2}+(m+q+1)u^m(u^2+b^2)^{q/2}.
\end{equation}

The logarithmic terms follow by integration by parts with
\begin{equation}
    \frac{\D}{\D u}\arsinh\frac{b}{u}=-\frac{b}{u\sqrt{u^2+b^2}};
\end{equation}
the boundary term $u^{m+1}\arsinh(b/u)$ vanishes as $u\to0^+$. \qed

\propTitle{Closed forms}{closedforms} For $p\in\{-1,1,3\}$, $m\in\{0,1,2,3\}$, $n\in\{0,1\}$ and $(a,b)\in Q^\circ$,
\begin{equation}
    \label{eq:Gform}
    \begin{split}
    \Gf{p}{m}{n}(a,b) &= \frac{1}{D}\Bigg(P_0 + P_1\sqrt{a^2+b^2} + P_2\log\left(\frac{\sqrt{a^2+b^2}+a}{b}\right)\\
    &\phantom{{}=\frac{1}{D}\Bigg(}+ P_3\log\left(\frac{\sqrt{a^2+b^2}+b}{a}\right)\Bigg),
    \end{split}
\end{equation}
with an integer $D$ and homogeneous integer polynomials $P_0,\dots,P_3$ in $a$ and $b$, listed in Appendix~\ref{app:antiderivatives}. These
are all the antiderivatives the moments need; any other one follows from the symmetry
$\Gf{p}{n}{m}(b,a)=\Gf{p}{m}{n}(a,b)$.

\vspace{\baselineskip}
\proof
The formulas are obtained by integrating first in $v$ and then in $u$ with Lemma~\lemmaRef{oned}; every step
produces only terms of the form~\eqref{eq:Gform}.

To verify them, note that each formula of Appendix~\ref{app:antiderivatives} is continuous on $Q$, vanishes on the axes and
satisfies $\partial_a\partial_bF=a^mb^n(a^2+b^2)^{p/2}$; these properties were checked by exact symbolic computation
(SymPy).
Hence $F=\Gf{p}{m}{n}$ by Proposition~\propRef{gproperties}(iii). All formulas also agree
with 30-digit adaptive quadrature of~\eqref{eq:G} within $2\cdot10^{-25}$. \qed

\vspace{\baselineskip}
\noindent\textbf{Example.} For $p=-1$, \eqref{eq:ibp} gives directly
\begin{equation}
    \Gf{-1}{0}{0}(a,b) = \int_0^a\arsinh\frac{b}{u}\D u
    = a\log\left(\frac{\sqrt{a^2+b^2}+b}{a}\right) + b\log\left(\frac{\sqrt{a^2+b^2}+a}{b}\right),
\end{equation}
the classical potential integral of a uniformly charged rectangle at its corner.
